\begin{lemma}
\label{lemma-relative-dualizing-if-have-omega}
Let $\varphi : A \to B$ be a finite type homomorphism of Noetherian rings.
Let $\omega_A^\bullet$ be a dualizing complex for $A$. Set
$\omega_B^\bullet = \varphi^!(\omega_A^\bullet)$. Denote
$D_A(K) = R\Hom_A(K, \omega_A^\bullet)$ for $K \in D_{\textit{Coh}}(A)$
and
$D_B(L) = R\Hom_B(L, \omega_B^\bullet)$ for $L \in D_{\textit{Coh}}(B)$.
Then there is a functorial isomorphism
$$
\varphi^!(K) = D_B(D_A(K) \otimes_A^\mathbf{L} B)
$$
for $K \in D_{\textit{Coh}}(A)$.
\end{lemma}

\begin{proof}
Observe that $\omega_B^\bullet$ is a dualizing complex for $B$ by
Lemma \ref{lemma-shriek-dualizing-algebraic}.
Let $A \to B \to C$ be finite type homomorphisms of Noetherian rings.
If the lemma holds for $A \to B$ and $B \to C$, then the lemma holds for
$A \to C$. This follows from
Lemma \ref{lemma-composition-shriek-algebraic}
and the fact that $D_B \circ D_B \cong \text{id}$ by
Lemma \ref{lemma-dualizing}.
Thus it suffices to prove the lemma in case $A \to B$ is
a surjection and in the case where $B$ is a
polynomial ring over $A$.

\medskip\noindent
Assume $B = A[x_1, \ldots, x_n]$. Since $D_A \circ D_A \cong \text{id}$,
it suffices to prove
$D_B(K \otimes_A B) \cong D_A(K) \otimes_A B[n]$ for $K$
in $D_{\textit{Coh}}(A)$.
Choose a bounded complex $I^\bullet$ of injectives representing
$\omega_A^\bullet$. Choose a quasi-isomorphism
$I^\bullet \otimes_A B \to J^\bullet$ where $J^\bullet$
is a bounded complex of $B$-modules. Given a complex
$K^\bullet$ of $A$-modules, consider the obvious
map of complexes
$$
\Hom^\bullet(K^\bullet, I^\bullet) \otimes_A B[n]
\longrightarrow
\Hom^\bullet(K^\bullet \otimes_A B, J^\bullet[n])
$$
The left hand side represents $D_A(K) \otimes_A B[n]$ and the right hand
side represents $D_B(K \otimes_A B)$. Thus it suffices to prove this
map is a quasi-isomorphism if the cohomology modules
of $K^\bullet$ are finite $A$-modules. Observe that the
cohomology of the complex in degree $r$ (on either side)
only depends on finitely many of the $K^i$. Thus we may
replace $K^\bullet$ by a truncation, i.e., we may assume
$K^\bullet$ represents an object of $D^-_{\textit{Coh}}(A)$.
Then $K^\bullet$ is quasi-isomorphic to a bounded
above complex of finite free $A$-modules.
Therefore we may assume $K^\bullet$ is a bounded
above complex of finite free $A$-modules.
In this case it is easy to that the
displayed map is an isomorphism of complexes which finishes
the proof in this case.

\medskip\noindent
Assume that $A \to B$ is surjective. Denote $i_* : D(B) \to D(A)$
the restriction functor and recall that $\varphi^!(-) = R\Hom(A, -)$
is a right adjoint to $i_*$ (Lemma \ref{lemma-right-adjoint}).
For $F \in D(B)$ we have
\begin{align*}
\Hom_B(F, D_B(D_A(K) \otimes_A^\mathbf{L} B))
& =
\Hom_B((D_A(K) \otimes_A^\mathbf{L} B) \otimes_B^\mathbf{L} F,
\omega_B^\bullet) \\
& =
\Hom_A(D_A(K) \otimes_A^\mathbf{L} i_*F, \omega_A^\bullet) \\
& =
\Hom_A(i_*F, D_A(D_A(K))) \\
& =
\Hom_A(i_*F, K) \\
& =
\Hom_B(F, \varphi^!(K))
\end{align*}
The first equality follows from More on Algebra, Lemma
\ref{more-algebra-lemma-internal-hom} and the definition
of $D_B$. The second equality by the adjointness mentioned
above and the equality
$i_*((D_A(K) \otimes_A^\mathbf{L} B) \otimes_B^\mathbf{L} F) =
D_A(K) \otimes_A^\mathbf{L} i_*F$
(More on Algebra, Lemma \ref{more-algebra-lemma-derived-base-change}).
The third equality follows from More on Algebra, Lemma
\ref{more-algebra-lemma-internal-hom}. The fourth because
$D_A \circ D_A = \text{id}$. The final equality by adjointness again.
Thus the result holds by the Yoneda lemma.
\end{proof}